% TOP_PROBLEM: {"id": 90, "title": "Hilbert's 17th Problem: Expression of Definite Forms", "category_id": 4, "category": {"id": 4, "name": "algebra", "display_name": "Algebra", "description": "Group theory, ring theory, field theory, and algebraic structures.", "slug": "algebra", "order_index": 4, "created_at": "2026-07-31T15:26:25.671Z"}, "status": "open"}
% FIRST_POSTED: null
% ENTRY_KIND: statement_only
% Imported from: batch02.tex; UnsolvedMath ID 90
\documentclass[11pt]{article}
\usepackage[margin=25mm]{geometry}
\usepackage{fontspec}
\setmainfont{Latin Modern Roman}
\usepackage{amsmath,amssymb,mathtools}
\usepackage{unicode-math}
\setmathfont{Latin Modern Math}
\usepackage{xurl,hyperref}
\hypersetup{colorlinks=true,urlcolor=blue}
\setcounter{secnumdepth}{0}
\setlength{\parindent}{0pt}
\setlength{\parskip}{5pt}
\setlength{\emergencystretch}{3em}
\newcommand{\sref}[2]{\hyperlink{source:#1:#2}{[S#2]}}
\newcommand{\cataloguescope}[1]{\paragraph{Recorded target.}#1}
\begin{document}
% UnsolvedMath ID 90; source catalogue code HIL-017
\setcounter{section}{89}
\section{Hilbert's 17th Problem: Expression of Definite Forms}\label{90}
{\small\sffamily UnsolvedMath ID 90 · Algebra}
\subsection{Definitions and mathematical statement}

\paragraph{Shared notation.}$\mathbb N=\{1,2,\ldots\}$, and $\mathbb Z,\mathbb Q,\mathbb R,\mathbb C$ denote the integers, rationals, reals and complex numbers. Any other conventions are specified locally.


For $n\ge1$, let $\mathbb R(x_1,\ldots,x_n)$ be the field of rational functions in $n$ independent variables with real coefficients. A rational function $f=p/q$, where $p,q\in\mathbb R[x_1,\ldots,x_n]$ and $q\ne0$, is \emph{nonnegative} if $p(a)/q(a)\ge0$ for every $a\in\mathbb R^n$ with $q(a)\ne0$. This condition is independent of the chosen representation. A \emph{sum of squares of rational functions} means a finite sum whose equality to $f$ holds in the rational-function field.
\paragraph{Statement recorded in the supplied source.}
For every $n\ge1$ and every nonnegative $f\in\mathbb R(x_1,\ldots,x_n)$, do there exist $m\ge1$ and $r_1,\ldots,r_m\in\mathbb R(x_1,\ldots,x_n)$ such that
\[
f=\sum_{j=1}^{m}r_j^2\quad\text{in }\mathbb R(x_1,\ldots,x_n)?
\]
\sref{90}{1}\sref{90}{2}
The polynomial formulation asks the same question for a polynomial nonnegative on all of $\mathbb R^n$. It is equivalent: if $f=p/q$ is nonnegative on $q\ne0$, then the polynomial $pq$ is nonnegative there, hence everywhere by continuity, and a rational-square representation of $pq$ can be divided by $q^2$. Artin answered this historical question affirmatively in 1927. \sref{90}{2}
\subsection{Short English statement}
Can every real rational function that is nonnegative wherever defined be written as a finite sum of squares of real rational functions? Artin's solution of Hilbert's seventeenth problem gives an affirmative answer.
\subsection{Sources}
\begin{description}
\item[\hypertarget{source:90:1}{S1}] ULAM.AI and the UnsolvedMath Contributors, \emph{UnsolvedMath: A Curated Collection of Open Mathematics Problems}, record 90 (HIL-017). CC BY 4.0. \url{https://huggingface.co/datasets/ulamai/UnsolvedMath}.
\item[\hypertarget{source:90:2}{S2}] Valéry Mahé, \emph{Using hyperelliptic curves to find positive polynomials that are not a sum of three squares in $\mathbb R(x,y)$}, arXiv:math/0703722v2, 13 September 2007, introduction, p. 2: definition of nonnegative polynomials, the rational-square formulation, and Artin's 1927 affirmative answer. \url{https://arxiv.org/pdf/math/0703722v2}.
\end{description}
\label{end:90}
\end{document}
